\begin{frame}{Positional Encoding: 순서 정보는 어떻게 넣나}
  어텐션 연산 자체는 \textbf{순서를 전혀 구분하지 않는다} —
  $QK^T$ 는 순서를 뒤섞어도 각 쌍의 관련도 자체는 똑같이
  계산(집합처럼 취급). 그런데 ``강아지가 고양이를 쫓는다''에서는
  \textbf{순서가 의미를 바꾼다}. 그래서 Transformer는 각 단어의
  임베딩에 위치 정보를 담은 벡터(positional encoding,
  \textbf{사인/코사인} 함수의 고정 패턴)를 \textbf{더해서}
  넣는다 — 같은 단어라도 \textbf{위치가 다르면 입력 벡터 자체가
  달라지}므로, 어텐션이 \textbf{간접적으로} 순서를 활용할 수
  있다.
  \vskip0.4em
  ``어텐션은 순서를 구분 못 한다''는 단순한 직관이 아니라,
  \textbf{코드로 확인할 수 있는 정리}:
  \begin{itemize}
    \item 단어 3개 $[\text{개}, \text{고양이}, \text{쫓는다}]$
      (임베딩 $(1,0), (0,1), (2,0)$)와 앞 두 단어를 바꾼
      $[\text{고양이}, \text{개}, \text{쫓는다}]$ 에 PE 없이
      self-attention($Q=K=V$) 적용.
    \item 두 문장의 어텐션 가중치 행렬:
      \[ W_{\text{B}} = P\, W_{\text{A}}\, P^T \]
      (완전히 같은 관계 — $P$ 는 1·2행을 바꾼
      \textbf{퍼뮤테이션 행렬}).
    \item 각 단어의 출력도 \textbf{정확히 퍼뮤테이션}: 문장 A에서
      ``고양이''의 출력 $[0.745,\; 0.503]$ , 문장 B에서
      ``고양이''의 출력도 \textbf{$[0.745,\; 0.503]$} —
      \emph{같은 단어의 표현이 위치에 무관하게 정확히
      같}다. ``누가 누구를 쫓는지''는 \textbf{통째로
      사라졌다}.
  \end{itemize}
  \vskip0.3em
  \begin{block}{permutation-equivariant(퍼뮤테이션-등변)}
    자기 어텐션 전체 연산이 \textbf{입력을 재배열하면 출력도
    같은 방식으로 재배열}. 순서 정보를 \emph{어딘가}에 덧붙여
    주지 않는 한 ``개가 고양이를 쫓는다''와 ``고양이가 개를
    쫓는다''를 구분하는 것은 \textbf{수학적으로 불가능}.
  \end{block}
\end{frame}
